\clearpage\begingroup
\HistoricalHead{160}{L}{PHYSIK.\quad BRIEFWECHSEL.}
\newcommand{\DioptrikCurrencyTwentyOne}{%
\begin{tikzpicture}[x=.13pt,y=-.13pt,baseline=-10pt]
\path[use as bounding box] (0,0) rectangle (75,89);
\path[fill=black,even odd rule] (58,9) -- (50,9) -- (49,10) -- (48,10) -- (45,14) -- (45,17) -- (44,18) -- (44,21) -- (42,23) -- (37,23) -- (36,24) -- (35,24) -- (31,28) -- (30,28) -- (30,29) -- (29,30) -- (29,34) -- (31,36) -- (31,40) -- (35,43) -- (36,43) -- (39,46) -- (42,46) -- (43,47) -- (43,53) -- (44,54) -- (44,55) -- (42,57) -- (38,57) -- (36,55) -- (35,55) -- (31,51) -- (30,51) -- (29,50) -- (29,49) -- (27,47) -- (27,46) -- (24,43) -- (21,43) -- (21,44) -- (20,45) -- (20,47) -- (19,48) -- (18,47) -- (18,45) -- (17,44) -- (14,43) -- (12,41) -- (10,43) -- (7,43) -- (5,46) -- (5,50) -- (4,51) -- (5,56) -- (8,56) -- (10,58) -- (9,59) -- (9,60) -- (11,63) -- (14,63) -- (15,62) -- (17,62) -- (17,61) -- (18,60) -- (20,63) -- (24,63) -- (26,61) -- (26,60) -- (22,56) -- (22,53) -- (24,51) -- (25,52) -- (28,53) -- (29,54) -- (30,57) -- (33,57) -- (36,60) -- (39,60) -- (40,61) -- (44,61) -- (45,62) -- (45,76) -- (41,80) -- (38,80) -- (37,79) -- (37,77) -- (36,76) -- (36,75) -- (33,72) -- (30,72) -- (29,73) -- (29,75) -- (27,77) -- (28,77) -- (29,78) -- (28,81) -- (31,84) -- (33,84) -- (34,85) -- (40,85) -- (42,83) -- (45,83) -- (46,82) -- (47,82) -- (48,81) -- (48,80) -- (51,77) -- (51,75) -- (53,73) -- (52,72) -- (52,69) -- (53,68) -- (53,64) -- (52,63) -- (52,62) -- (54,60) -- (55,61) -- (56,60) -- (59,60) -- (63,57) -- (66,57) -- (67,56) -- (68,56) -- (68,54) -- (69,53) -- (70,53) -- (70,51) -- (72,49) -- (72,44) -- (71,43) -- (72,42) -- (72,36) -- (70,34) -- (71,33) -- (66,29) -- (64,29) -- (63,28) -- (63,27) -- (65,25) -- (65,21) -- (63,19) -- (63,16) -- (62,15) -- (61,12) -- cycle
(10,48) -- (11,47) -- (14,47) -- (15,48) -- (15,52) -- (14,53) -- (12,53) -- (11,54) -- (10,54) -- (9,53) -- (9,52) -- (10,51) -- cycle
(65,35) -- (69,39) -- (69,44) -- (68,45) -- (68,47) -- (63,51) -- (63,53) -- (62,54) -- (57,55) -- (53,58) -- (52,57) -- (52,45) -- (53,44) -- (55,44) -- (56,43) -- (57,43) -- (62,38) -- (62,37) -- (64,35) -- cycle
(51,26) -- (52,26) -- (53,27) -- (56,27) -- (61,30) -- (61,31) -- (60,32) -- (61,33) -- (60,34) -- (59,37) -- (56,40) -- (54,40) -- (53,41) -- (52,41) -- (51,40) -- (51,28) -- (50,27) -- cycle
(41,26) -- (42,26) -- (43,27) -- (43,31) -- (41,33) -- (42,33) -- (44,35) -- (44,37) -- (43,38) -- (44,39) -- (44,42) -- (43,43) -- (41,41) -- (36,40) -- (33,37) -- (33,34) -- (34,33) -- (34,32) -- (39,27) -- (40,27) -- cycle
(53,14) -- (57,14) -- (58,15) -- (58,17) -- (61,21) -- (61,22) -- (60,23) -- (61,24) -- (60,25) -- (58,25) -- (57,24) -- (53,24) -- (52,23) -- (51,23) -- (48,20) -- cycle;
\end{tikzpicture}}
\newcommand{\DioptrikMoonTwentyOne}{%
\begin{tikzpicture}[x=.12pt,y=-.12pt,baseline=-8pt,line width=.5pt]
\path[use as bounding box] (0,0) rectangle (60,76);
\draw (54,5) .. controls (3,2) and (3,72) .. (54,71)
 .. controls (20,61) and (20,15) .. (54,5);
\end{tikzpicture}}

\begin{center}
[18.]\par\medskip
\textbf{\textsc{Encke} an \textsc{Gauss}.\quad Berlin, 30. October 1840.}
\end{center}
\noindent \{. . . . . . . . . . . . . . . . . . . . .\par
\noindent\hspace*{\parindent}Über das dialytische Princip kenne ich nur eine Abhandlung von \textsc{Rogers} in den Memoirs der Astronomical Society of London, Vol. III, pag. 229: \textit{On}
\par\clearpage\SourceLeaf{166}\HistoricalHead{161}{R}{GAUSS AN ENCKE.}
\noindent\textit{the construction of large achromatic telescopes}, by \textsc{Alexander Rogers}; Read Apr. 11. 1828. --- Es ist diese Abhandlung, welche eine Prioritätsfrage zwischen \textsc{Rogers} und \textsc{Plössl} hervorgerufen hat, welche wie ich glaube zu Gunsten des erstern entschieden werden muss. Früher hatte ich gesucht, mich etwas mit der Theorie bekannt zu machen, und finde als Notiz bei \textsc{Rogers} Vorschlag beigeschrieben, dass entweder die Winkel, unter welchen die verschiedenfarbigen Strahlen zusammen kommen, oder die Punkte, in welchen sie die Axe schneiden, weiter von einander verschieden sind, als dass man ein gutes Resultat erwarten könnte. \textsc{Plössl} schiene mir durch mechanische Nachhilfe und Probiren der besten Entfernung, in welcher jedes Okular von der zusammengesetzten Linse zu stellen seyn mögte, diese Nachtheile praktisch so viel als möglich ausgeglichen zu haben. Es hat mich ungemein gefreut, dass Ihre Untersuchung ebenfalls Nachtheile der Construktion nachgewiesen hat. . . . . . . .\}
\begin{center}
[19.]\par\medskip
\textbf{\textsc{Gauss} an \textsc{Encke}.\quad Göttingen, 23. Dezember 1840.}
\end{center}
\noindent\hspace*{\parindent}Indem ich Ihnen, hochgeschätztester Freund, hieneben das dialytische Fernrohr mit dem verpflichtetsten Danke zurückschicke, muss ich zugleich um Verzeihung bitten, dass ich es so lange behalten. Schuld hat der Umstand, dass fast gleichzeitig mit dem Empfange ich von \textsc{Schumacher} die Nachricht erhielt, dass er ein ganz ähnliches (neues, schon vor 4 Jahren bestellt gewesenes aber nun disponibel gewordenes) an mich abgeschickt habe. Da sich dadurch die angenehme Gelegenheit eröffnete, zwei Instrumente neben einander zu vergleichen, und ich das Instrument täglich erwarten durfte, so ging in vergeblichem Warten ein Tag nach dem andern hin. Hinterdrein hat sich nun gefunden, dass es erst lange bei dem Spediteur in Hamburg liegen geblieben ist und auf dem Transport durch Fracht fast 3 Wochen zugebracht hat; erst vor 8 Tagen ist es angekommen, und anfangs war es wegen mehrfach mangelhafter mechanischer Arbeit gar nicht zu gebrauchen. Letzteres ist nun redressirt, und so habe ich einige Vergleichungen, freilich nur auf der Erde, angestellt. Die Wirkung der terrestrischen Oculare zeigt sich an beiden Instrumenten so gleich, dass ich keinen Unterschied angeben kann. Die
\par\HistoricalFoot{XI\textsubscript{1}.}{21}
\clearpage\SourceLeaf{167}\HistoricalHead{162}{L}{PHYSIK.\quad BRIEFWECHSEL.}
\noindent stärksten astronomischen Oculare hingegen zeigen einen, obwohl nicht grossen Unterschied zum Vortheil Ihres Instruments. Allein die stärksten Vergrösserungen sind nicht bei beiden Instrumenten dieselben; an dem Ihrigen hatte ich sie früher $=99$ gemessen (die schwächere 62); an dem andern habe ich sie noch nicht bestimmt, allein sie ist augenscheinlich kleiner, vielleicht so wie in \textsc{Plössls} Catalog angegeben wird, 85. Da die Gewinde der Oculareinsätze bei beiden Instrumenten beinahe gleich sind, so liess sich das Ocular des Ihrigen bequem in das andere Instrument einschrauben, und dann war nach meinem Schätzen die Wirkung ganz dieselbe, wie an dem Ihrigen. Ich habe die Absicht, jenes Instrument zu behalten, und ich denke dann künftig noch ein neues Ocular dazu machen zu lassen, da, so viel ich bis jetzt nach der Theorie beurtheilen kann, die \textsc{Plössl}schen nicht dieser gemäss combinirt sind. Ich meine, dass wenigstens für die schwächere Vergrösserung durch eine andere Ocularanordnung \textls[120]{ohne sonstigen Nachtheil} ein grösseres Gesichtsfeld sich erzielen lassen wird.

Was ich Ihnen früher geschrieben habe, nemlich dass das [durch] dialytisch angeordnete Objectivlinsen hervorgebrachte Bild farbenrein \textls[120]{nicht sein kann}, bleibt unumstösslich wahr; allein eine tiefer eindringende Theorie hat mich überzeugt, dass dieses durch die Oculargläser wieder aufgehoben werden [kann] und, worauf es doch am Ende allein ankommt, \textls[120]{im Auge} ein vollkommen farbenreines Bild möglich ist, während diess bei einem Objectiv der gewöhnlichen Einrichtung theoretisch unmöglich ist, wenigstens nicht möglich für mehr als Einen Oculareinsatz. Ich habe einige Bemerkungen darüber einer Abhandlung beigefügt, die der Societät übergeben und jetzt bereits unter der Presse ist[*)];\begingroup\renewcommand{\thefootnote}{}\footnotetext{[*) \textit{Dioptrische Untersuchungen}, Werke V, S. 243, insbesondere S. 274, 275.]}\endgroup\ schon 2 Bogen [sind] gedruckt. Der Inhalt dieser Abhandlung bezieht sich auf eine mir seit 40---45 Jahren eigenthümliche Weise, die Elemente der Dioptrik zu behandeln, und wovon etwas zu publiciren ich bisher nur darum Abstand genommen hatte, weil mir der Gegenstand gar zu elementarisch schien. Indessen dürfte sie doch wohl nicht überflüssig sein, da die gewöhnliche Behandlungsweise dioptrischer Rechnungen einen Charakter von Ungenauigkeit, ich möchte sagen Rohheit hat, der den mathematischen Sinn sehr beleidigt. \textsc{Bessel} ist, deucht mir, der erste, der in dioptrische Versuche
\par\clearpage\SourceLeaf{168}\HistoricalHead{163}{R}{GAUSS, ENCKE, GERLING.}
\noindent über schon fertige Linsengläser eine anständige Präcision zu bringen gesucht hat[*)];\begingroup\renewcommand{\thefootnote}{}\footnotetext{[*) F. W. \textsc{Bessel}, \textit{Ueber ein Mittel zur Bestimmung der Brennweite des Objectivglases eines Fernrohrs}, Astronomische Nachrichten \textbf{17}, 1840, Spalte 289 und \textit{Grundformeln der Dioptrik}, ebenda \textbf{18}, 1841, Spalte 97; Abhandlungen herausg. von \textsc{Engelmann} II, S. 107 und III, S. 282.}\endgroup\ nur muss man sich wundern, wie er in den Irrthum gerathen ist, als seien die Erscheinungen bei einem Objectiv von endlicher Dicke genau dieselben, wie sie bei einer im optischen Mittelpunkte des erstern befindlichen Linse von unendlich kleiner Dicke, von derselben Brennweite sein würden. In folge dieses Irrthums hat er eine ganz falsche Brennweite herausgebracht und ich würde mich nicht wundern, wenn der Fehler anstatt (nach seiner Schätzung) $\frac{1}{75000}$ zu sein, $\frac{1}{1500}$ betrüge. Ich muss Sie übrigens bitten, diese Mittheilung nur wie eine \textls[120]{bloss Ihnen vertraulich} gemachte zu betrachten.
\par\noindent . . . . . . . . . . . . . . . . . . . . .
\begin{center}
[20.]\par\medskip
\textbf{\textsc{Encke} an \textsc{Gauss}.\quad Berlin, 31. Dezember 1840.}
\end{center}
\noindent \{. . . . . . . . . . . . . . . . . . . . .\par
\noindent\hspace*{\parindent}Ihrer neuen Abhandlung über das dialytische Fernrohr sehe ich mit umso grösserem Verlangen entgegen, als Sie schon 1812 Ihre Unzufriedenheit mit der Behandlung solcher Aufgaben mir geäussert haben und ich so hoffen darf, dass auch auf diesem Felde wie auf jedem, welches Sie betreten, der neue und wahre Weg gezeigt werden wird. . . . . . . .\}
\begin{center}
[21.]\par\medskip
\textbf{\textsc{Gauss} an \textsc{Gerling}.\quad Göttingen, 12. Mai 1841.}
\end{center}
\noindent . . . . . . . . . . . . . . . . . . . . .\par
\noindent\hspace*{\parindent}Ich habe vorigen Winter ein kleines dialytisches Fernrohr angeschafft (Kosten cca. 100\,\DioptrikCurrencyTwentyOne) dessen leichte Handhabung und in seiner Art schon reines Bild mir bei \DioptrikMoonTwentyOne\ Betrachtungen viel Genuss macht. . . . . . . .
\clearpage\endgroup
